% !TEX root = ../main.tex
\chapter{大数定律}
\label{ch:lln}

\noindent\emph{用到的前置工具：随机变量、期望与方差（概率与测度部分）；收敛模式——依概率收敛 $\pto$、几乎必然收敛 $\asto$（见前一章``收敛模式'')；Lebesgue 积分意义下的期望与基本不等式（概率与测度部分）。}
\medskip

经济学家手里几乎所有的``用数据说话'',归根结底都依赖同一句承诺：\emph{把样本平均起来，就能逼近你真正关心的那个总体量}。问卷里 1000 个人的平均收入，能代表全人群的期望收入吗？把过去 240 个月的股票超额收益平均一下，能当作真实的风险溢价吗？蒙特卡洛里抽 $10^6$ 个随机数取平均，能算出那个解析上写不出来的积分吗？这些操作天天在做，但它们\emph{凭什么}成立？大数定律（Law of Large Numbers，LLN）正是这句承诺的严格兑现：在适当条件下，样本均值 $\bar X_n$ 会收敛到总体期望 $\mu=\E{X}$。

本章要做三件事。先点明这条定律为何是整个统计推断的地基——没有它，``用样本估计总体''就成了无据的信仰；再把它\emph{证出来}，从一条朴素得近乎显然的不等式（Markov 不等式）出发，三步到达弱大数定律，并交代去掉方差假设需要的``截断''技巧；最后说清它通向何处——矩估计、GMM、极大似然的一致性，蒙特卡洛积分，经验分布函数，乃至时间序列里的遍历定理，骨子里都是大数定律在各处的化身。

\section{它在守护什么：从样本到总体的桥}
\label{sec:lln-motivation}

把问题摆清楚。我们关心一个总体的某个数字特征——通常是某随机变量 $X$ 的期望 $\mu=\E{X}$（``平均收入''``平均回报''``某事件的真实概率''都属此类，后者是 $X=\indicator_A$ 的期望）。我们手上没有整个总体，只有一组观测 $X_1,\dots,X_n$。最自然的估计就是\textbf{样本均值}
\[
    \bar X_n=\frac1n\sum_{i=1}^n X_i .
\]
问题是：凭什么相信这个由\emph{有限个、带随机性}的观测算出来的数，会接近那个我们看不见的 $\mu$？

直觉上的理由人人都有：样本越大，``运气''越被摊薄。但直觉不是定理。$\bar X_n$ 本身是随机变量——换一组样本就换一个值，它永远不会\emph{恰好}等于 $\mu$。所谓``收敛'',必须是某种概率意义下的收敛（这正是前一章``收敛模式''要分清依概率 $\pto$、几乎必然 $\asto$、依分布 $\dto$ 的原因）。大数定律的内容，就是把``样本越大越准''这句口语，钉成一个可证的极限命题：

\state{大数定律是一切``样本矩估计总体矩''的执照}{
几乎每一种估计方法的一致性，最终都化约为一句``某个样本平均收敛到对应的总体期望''：矩估计直接拿样本矩当总体矩，GMM 让样本矩条件趋于零，极大似然让样本平均对数似然趋于其期望，蒙特卡洛拿样本平均当积分。\emph{大数定律就是这张执照}——没有它，``$n$ 大了就准''只是信念；有了它，它是定理。后续``极值估计的统一框架''一章会反复调用本章的结论。
}

\section{两件趁手的不等式}
\label{sec:lln-inequalities}

证明大数定律出乎意料地省力，靠的是一条几乎显然、却异常有用的不等式。它把``期望''这一个数，翻译成对``尾部概率''的控制——而``$\bar X_n$ 偏离 $\mu$ 太远''恰恰就是一个尾部事件。

\thm{Markov 不等式}{
设 $Y\ge 0$ 是非负随机变量，$a>0$。则
\[
    \Prob{Y\ge a}\;\le\;\frac{\E{Y}}{a} .
\]
}
\pf{
关键是一个``指示函数被放大''的逐点不等式。对任意结果 $\omega$，比较两个非负量：若 $Y(\omega)\ge a$，则 $\indicator_{\set{Y\ge a}}(\omega)=1$，而 $Y(\omega)/a\ge 1$，故 $\indicator_{\set{Y\ge a}}\le Y/a$；若 $Y(\omega)<a$，则左边为 $0$，而右边非负，不等式同样成立。于是逐点有
\[
    \indicator_{\set{Y\ge a}}\;\le\;\frac{Y}{a} .
\]
两边取期望（期望保序，且 $\E{\indicator_A}=\Prob{A}$）：
\[
    \Prob{Y\ge a}=\E{\indicator_{\set{Y\ge a}}}\le\E{\frac{Y}{a}}=\frac{\E{Y}}{a} . \qedhere
\]
}

\rmkb[一行证明里藏着全部直觉]{
Markov 不等式说的是一件朴素的事：\emph{一个非负量若平均值不大，它就不可能经常取很大的值}。若平均工资是 $5$ 万，那么挣到 $50$ 万以上的人占比至多 $1/10$——否则光这一撮人就把平均拉过 $5$ 万了。它把对\emph{均值}的一点知识，兑换成对\emph{尾部}的控制。整个大数定律的证明，就是把``$\bar X_n$ 远离 $\mu$''这件事识别成一个尾部事件，再用它来夹。
}

把 Markov 用到``偏离均值的平方''上，立刻得到一条以方差控制离散程度的不等式——它是弱大数定律的直接引擎。

\thm{Chebyshev 不等式}{
设随机变量 $X$ 有有限均值 $\mu=\E{X}$ 与有限方差 $\sigma^2=\Var{X}$。则对任意 $\ve>0$，
\[
    \Prob{\abs{X-\mu}\ge\ve}\;\le\;\frac{\sigma^2}{\ve^2} .
\]
}
\pf{
令 $Y=(X-\mu)^2\ge 0$，这是个非负随机变量，其期望恰是方差：$\E{Y}=\E{(X-\mu)^2}=\sigma^2$。注意事件等价
\[
    \set{\abs{X-\mu}\ge\ve}=\set{(X-\mu)^2\ge\ve^2}=\set{Y\ge\ve^2} ,
\]
对 $Y$ 用 Markov 不等式（取 $a=\ve^2$）：
\[
    \Prob{\abs{X-\mu}\ge\ve}=\Prob{Y\ge\ve^2}\le\frac{\E{Y}}{\ve^2}=\frac{\sigma^2}{\ve^2} . \qedhere
\]
}

Chebyshev 不等式的形态值得记牢：\textbf{偏离均值至少 $\ve$ 的概率，被方差与 $\ve^2$ 的比值压住}。方差越小、$X$ 越集中，这个上界越紧。大数定律的诀窍正在于：样本均值 $\bar X_n$ 的方差会随 $n$ 趋于零，于是把 $X$ 换成 $\bar X_n$、把这条上界往零里赶，就得到了收敛。

\section{弱大数定律：三步证完}
\label{sec:lln-wlln}

先在最透亮的假设下把定理拿下：观测两两不相关、同均值、且有相同的有限方差。这覆盖了 i.i.d.（独立同分布）有限方差这一最常用的情形，而证明短到可以一眼看穿。

\asm{两两不相关、同均值、方差有界}{
$X_1,X_2,\dots$ 各有相同均值 $\E{X_i}=\mu$ 与相同方差 $\Var{X_i}=\sigma^2<\infty$，且两两不相关：$i\neq j$ 时 $\Cov{X_i}{X_j}=0$。（i.i.d.\ 是其特例。）
}

\thm{弱大数定律（WLLN，有限方差）}{
在上述假设下，样本均值\emph{依概率}收敛到总体均值：
\[
    \bar X_n=\frac1n\sum_{i=1}^n X_i\;\pto\;\mu ,
\]
即对任意 $\ve>0$，$\ \Prob{\abs{\bar X_n-\mu}\ge\ve}\to 0\ (n\to\infty)$。
}
\pf{
三步：算 $\bar X_n$ 的均值、算它的方差、套 Chebyshev。

\emph{第一步：均值不变。} 期望是线性的，
\[
    \E{\bar X_n}=\frac1n\sum_{i=1}^n\E{X_i}=\frac1n\cdot n\mu=\mu .
\]
所以 $\bar X_n$ 是 $\mu$ 的\emph{无偏}估计——它没有系统性的偏差，平均而言瞄准了靶心。

\emph{第二步：方差像 $1/n$ 那样塌缩。} 这是全证明的心脏。由两两不相关，和的方差等于方差之和（交叉的协方差项全为零）：
\[
    \Var{\bar X_n}
    =\Var{\frac1n\sum_{i=1}^n X_i}
    =\frac{1}{n^2}\sum_{i=1}^n\Var{X_i}
    =\frac{1}{n^2}\cdot n\sigma^2
    =\frac{\sigma^2}{n}\;\xrightarrow[n\to\infty]{}\;0 .
\]
注意那个 $1/n^2$：把和缩成均值时除了一个 $n$，方差里就变成除以 $n^2$，而被加的项只有 $n$ 个——一进一出，净剩 $1/n$。\emph{这就是``样本越大越准''的全部数学内容}：估计量绕着靶心的抖动，按 $1/n$ 的速度被摊平。

\emph{第三步：Chebyshev 收口。} 把 Chebyshev 不等式用到 $\bar X_n$ 上（它的均值是 $\mu$、方差是 $\sigma^2/n$）：对任意 $\ve>0$，
\[
    \Prob{\abs{\bar X_n-\mu}\ge\ve}\;\le\;\frac{\Var{\bar X_n}}{\ve^2}=\frac{\sigma^2}{n\,\ve^2}\;\xrightarrow[n\to\infty]{}\;0 .
\]
固定 $\ve$ 让 $n\to\infty$，右边趋于零，左边被夹到零。这正是 $\bar X_n\pto\mu$ 的定义。\qedhere
}

\rmkb[为什么叫``弱'']{
这里证出的是\emph{依概率}收敛：对每个固定的 $n$，``$\bar X_n$ 偏离 $\mu$ 超过 $\ve$''的概率很小，且随 $n$ 趋于零。但它\emph{没有}断言：沿着\emph{某一条具体的样本路径} $\bar X_1,\bar X_2,\dots$ 看过去，这串数最终会稳稳停在 $\mu$ 附近不再大幅跳出。后者是更强的``几乎必然''收敛，由强大数定律给出。``弱''与``强''的分野，正是前一章``收敛模式''里 $\pto$ 与 $\asto$ 的分野——记住：$\asto\Rightarrow\pto$，反之不然。
}

图~\ref{fig:lln-1} 把第二步的直觉画了出来：样本均值的分布随 $n$ 增大\emph{越挤越窄}地堆到 $\mu$ 上，于是落在 $\mu$ 的任意 $\ve$-邻域之外的概率被挤向零。

\begin{figure}[h]
\centering
\begin{tikzpicture}[scale=1.0,>=Stealth]
    % 坐标轴：x 表示 Xbar 的取值，y 表示密度
    \draw[->] (0,0) -- (11.6,0) node[right] {$\bar X_n$ 的取值};
    \draw[->] (0,0) -- (0,5.0) node[above] {密度};
    % mu 处竖直参考线（虚线），x=5.5 对应取值 2
    \draw[dashed,gray] (5.5,0) -- (5.5,4.6);
    \node[below] at (5.5,0) {$\mu$};
    % epsilon-邻域：mu +- 0.3 -> X 4.4 与 6.6（用浅色竖线标出边界）
    \draw[densely dotted,gray!70] (4.4,0) -- (4.4,2.0);
    \draw[densely dotted,gray!70] (6.6,0) -- (6.6,2.0);
    \node[gray!60!black,anchor=north,font=\small] at (4.4,-0.05) {$\mu-\ve$};
    \node[gray!60!black,anchor=north,font=\small] at (6.6,-0.05) {$\mu+\ve$};
    % 三条密度曲线：Xbar_n ~ N(2, sigma^2/n)，sigma=1，n=4,16,64
    % 取值 v 映射 x=11*(v-0.5)/3；密度 d 映射 y=4.2*d/maxpeak, maxpeak=3.192(n=64)
    % 直接用高斯函数绘制，自变量 \v 走取值域 [0.5,3.5]
    % n=4: sd=0.5, peak=0.798 -> y_peak=1.050
    \draw[thick,blue!35!black,domain=0.5:3.5,smooth,samples=120,variable=\v]
        plot ({11*(\v-0.5)/3},{4.2/3.192*(1/(0.5*2.5066))*exp(-0.5*((\v-2)/0.5)^2)});
    % n=16: sd=0.25, peak=1.596 -> y_peak=2.100
    \draw[thick,blue!60!cyan,domain=0.5:3.5,smooth,samples=160,variable=\v]
        plot ({11*(\v-0.5)/3},{4.2/3.192*(1/(0.25*2.5066))*exp(-0.5*((\v-2)/0.25)^2)});
    % n=64: sd=0.125, peak=3.192 -> y_peak=4.200
    \draw[thick,cyan!70!black,domain=0.5:3.5,smooth,samples=220,variable=\v]
        plot ({11*(\v-0.5)/3},{4.2/3.192*(1/(0.125*2.5066))*exp(-0.5*((\v-2)/0.125)^2)});
    % 曲线标签：放在各自峰顶的右上空白，细引线指向曲线，不穿其它曲线
    \node[blue!35!black,anchor=west,font=\small] at (7.7,1.15) {$n=4$};
    \draw[blue!35!black,->,thin] (7.65,1.15) -- (6.55,1.0);
    \node[blue!60!cyan,anchor=west,font=\small] at (7.0,2.25) {$n=16$};
    \draw[blue!60!cyan,->,thin] (6.95,2.25) -- (6.05,2.05);
    \node[cyan!70!black,anchor=west,font=\small] at (6.2,4.3) {$n=64$};
    \draw[cyan!70!black,->,thin] (6.15,4.3) -- (5.65,4.1);
\end{tikzpicture}
\caption{样本均值 $\bar X_n$ 的分布随 $n$ 增大越挤越窄地集中到 $\mu$（这里 $\sigma=1$，方差 $\sigma^2/n$ 依次为 $1/4,1/16,1/64$）。落在 $\ve$-邻域 $[\mu-\ve,\mu+\ve]$ 之外的概率随之被挤向零，这正是 $\bar X_n\pto\mu$。}
\label{fig:lln-1}
\end{figure}

\section{去掉方差假设：Khinchin 与截断的骨架}
\label{sec:lln-khinchin}

Chebyshev 证法漂亮，但它欠下一笔债：用到了\emph{有限方差} $\sigma^2<\infty$。可经济学里偏偏有不少厚尾分布（收入、城市规模、金融收益）方差是发散的；对它们，``期望''仍然有限、有意义，我们照样想要``样本均值收敛到期望''。Khinchin 的定理正补上这一块：只要 i.i.d.\ 且\emph{期望有限}，弱大数定律就成立——方差可以不存在。

\thm{Khinchin 弱大数定律}{
设 $X_1,X_2,\dots$ 独立同分布，期望有限：$\E{\abs{X_1}}<\infty$，记 $\mu=\E{X_1}$。则
\[
    \bar X_n\;\pto\;\mu .
\]
}

完整证明（特征函数法只需两行，但要用到尚未介绍的特征函数）超出本章主线。这里给\emph{截断法}的骨架——它是渐近统计里反复出现的标准招式，值得看清其轮廓。困难在于方差可能无穷，Chebyshev 用不上；办法是把每个变量在一个随 $n$ 增长的阈值处``砍头'',分成``截断部分''与``尾巴部分''分别对付。

\state{截断法三步骨架（直觉）}{
取一个随 $n$ 缓慢增长的阈值 $b_n\to\infty$，把每个 $X_i$ 拆成截断与残余两块：
\[
    X_i=\underbrace{X_i\,\indicator_{\set{\abs{X_i}\le b_n}}}_{Y_i\text{（截断，有界）}}+\underbrace{X_i\,\indicator_{\set{\abs{X_i}> b_n}}}_{Z_i\text{（残余，少见）}} .
\]
\textbf{(1) 截断块 $Y_i$ 有界}，故方差有限，可对 $\bar Y_n$ 用 Chebyshev——这块照搬上一节的论证；只需把 $b_n$ 调得让方差项仍趋于零。\textbf{(2) 残余块可忽略}：因 $\E{\abs{X_1}}<\infty$，落在阈值之外的概率与贡献都随 $b_n\to\infty$ 被 $\E{\abs{X_1}\indicator_{\set{\abs{X_1}>b_n}}}\to 0$ 压住（这一步正是``期望有限''发挥作用之处）。\textbf{(3) 截断后的均值不偏太多}：$\E{Y_i}\to\mu$。把三块拼起，便得 $\bar X_n\pto\mu$。\emph{有限期望——而非有限方差——才是弱大数定律真正需要的最低要求}。
}

\rmkb[为什么期望有限不可再省]{
``期望有限''是底线，不能再弱。经典反例是柯西分布（密度 $\frac{1}{\pi(1+x^2)}$）：它对称、看着``中心在 0'',但 $\E{\abs{X}}=\infty$，期望根本不存在。对它，$\bar X_n$ \emph{不会}收敛到任何常数——事实上 $\bar X_n$ 与单个 $X_1$ 同分布，再多的样本也不会让平均值安定下来。这正说明：大数定律不是``取平均''这一动作本身的魔法，而是``存在一个有限期望可供收敛''这一前提的兑现。}

\section{强大数定律：沿路径的收敛}
\label{sec:lln-slln}

弱大数定律谈的是``每个 $n$ 上的概率'';强大数定律谈的是``整条样本路径的命运''。把同一个 $\omega$（一整次实验、一整条无穷长的观测序列）固定下来，就得到一串确定的数 $\bar X_1(\omega),\bar X_2(\omega),\dots$。强大数定律断言：\emph{除了一个概率为零的坏样本集合之外，每一条这样的路径作为普通数列都收敛到 $\mu$}。

\thm{Kolmogorov 强大数定律（SLLN）}{
设 $X_1,X_2,\dots$ 独立同分布且期望有限 $\E{\abs{X_1}}<\infty$，$\mu=\E{X_1}$。则样本均值\emph{几乎必然}收敛到 $\mu$：
\[
    \Prob{\,\lim_{n\to\infty}\bar X_n=\mu\,}=1 ,\qquad\text{即}\quad \bar X_n\;\asto\;\mu .
\]
（反之亦真：若 $\bar X_n$ 几乎必然收敛到某有限常数，则 $\E{\abs{X_1}}$ 必有限。故对 i.i.d.\ 而言，有限期望既充分又必要。）
}

严格证明依赖 Kolmogorov 的极大不等式与三级数定理（或 Etemadi 的子序列截断论证），属于较重的纯概率构造，本书按使用者的需要只陈述、给直觉，不逐行搭建。要抓住的是``强''比``弱''强在哪里：

\state{$\asto$ 与 $\pto$ 的差别，画在一张图上才看得清}{
\emph{依概率}（弱）：盯住某个固定的时刻 $n$，问``此刻偏离超 $\ve$ 的概率多大''——它随 $n$ 趋零。但它\emph{允许}某条路径在 $n=10^3,10^6,10^9,\dots$ 处一次次偶发地跳出 $\ve$-带，只要这种跳出越来越稀疏即可。\\
\emph{几乎必然}（强）：盯住一整条路径，问``它最终会不会\emph{永远}留在 $\mu$ 的任意 $\ve$-带里''——强大数定律说，除一个零概率的例外集外，\emph{会}。也就是：对几乎每条路径，存在一个时刻 $N(\omega)$，此后再不跳出。\\
一句话：弱定律管``某一刻''，强定律管``从某刻起的全部将来''。$\asto\Rightarrow\pto$；反向一般不成立。}

图~\ref{fig:lln-2} 画几条真实模拟出来的样本均值路径：横轴 $n$（对数刻度，好让早期的剧烈摆动与后期的安定同框可见），纵轴 $\bar X_n$。早期各路径在宽处乱晃，随 $n$ 增大被一条按 $\propto 1/\sqrt n$ 收紧的波动带逼回，最终一起贴向水平线 $\mu$。强大数定律保证的，正是这种``逐条路径收敛''。

\begin{figure}[h]
\centering
\begin{tikzpicture}[scale=1.0,>=Stealth]
    % 画布：x 对数刻度 n in [1,1000] -> [0,11]；y 取值 v in [0.8,3.2] -> [0,5]
    \draw[->] (0,0) -- (11.8,0) node[right] {$n$（对数刻度）};
    \draw[->] (0,0) -- (0,5.3) node[above] {$\bar X_n$};
    % x 轴刻度：n=1,10,100,1000 -> X=0,3.667,7.333,11
    \foreach \x/\lab in {0/1,3.667/10,7.333/100,11/1000}{
        \draw (\x,0.06) -- (\x,-0.06) node[below,font=\small] {\lab};
    }
    % y 轴刻度：v=1.0,1.5,2.0,2.5,3.0 -> Y=0.417,1.458,2.5,3.542,4.583
    \foreach \y/\lab in {0.417/1.0,1.458/1.5,3.542/2.5,4.583/3.0}{
        \draw (0.06,\y) -- (-0.06,\y) node[left,font=\small] {\lab};
    }
    % 收敛波动带 mu +- 1.5 sigma/sqrt(n)，淡灰填充（先画，垫底）
    \fill[gray!12]
        (1.104,4.541) -- (1.749,4.167) -- (2.563,3.791) -- (3.311,3.521) -- (4.084,3.301)
        -- (4.848,3.130) -- (5.615,2.995) -- (6.381,2.889) -- (7.148,2.806) -- (7.914,2.741)
        -- (8.680,2.689) -- (9.447,2.649) -- (10.213,2.617) -- (11.000,2.591)
        -- (11.000,2.409) -- (10.213,2.383) -- (9.447,2.351) -- (8.680,2.311) -- (7.914,2.259)
        -- (7.148,2.194) -- (6.381,2.111) -- (5.615,2.005) -- (4.848,1.870) -- (4.084,1.699)
        -- (3.311,1.479) -- (2.563,1.209) -- (1.749,0.833) -- (1.104,0.459) -- cycle;
    % mu 水平线（粗，置于带之上、路径之下）
    \draw[thick,red!60!black] (0,2.5) -- (11.2,2.5);
    \node[red!60!black,anchor=west,font=\small] at (11.05,2.5) {$\mu$};
    % 四条样本路径（seed 21，坐标精确算自模拟）
    \draw[blue!35!black,thick]
        (1.104,2.016) -- (1.749,2.477) -- (2.563,2.080) -- (3.311,2.748) -- (4.084,2.243)
        -- (4.848,2.335) -- (5.615,2.289) -- (6.381,2.512) -- (7.148,2.534) -- (7.914,2.628)
        -- (8.680,2.516) -- (9.447,2.502) -- (10.213,2.456) -- (11.000,2.525);
    \draw[cyan!55!black,thick]
        (1.104,3.361) -- (1.749,2.520) -- (2.563,2.671) -- (3.311,2.824) -- (4.084,2.514)
        -- (4.848,2.415) -- (5.615,2.501) -- (6.381,2.444) -- (7.148,2.458) -- (7.914,2.561)
        -- (8.680,2.556) -- (9.447,2.522) -- (10.213,2.528) -- (11.000,2.500);
    \draw[violet!60!black,thick]
        (1.104,0.642) -- (1.749,0.601) -- (2.563,1.236) -- (3.311,2.140) -- (4.084,2.225)
        -- (4.848,2.490) -- (5.615,2.798) -- (6.381,2.947) -- (7.148,2.747) -- (7.914,2.681)
        -- (8.680,2.674) -- (9.447,2.600) -- (10.213,2.560) -- (11.000,2.487);
    \draw[teal!60!black,thick]
        (1.104,2.307) -- (1.749,2.181) -- (2.563,1.165) -- (3.311,2.375) -- (4.084,2.412)
        -- (4.848,2.461) -- (5.615,2.734) -- (6.381,2.662) -- (7.148,2.669) -- (7.914,2.638)
        -- (8.680,2.537) -- (9.447,2.571) -- (10.213,2.579) -- (11.000,2.478);
    % 波动带标签：停在左上大片空白，细引线指向带的上沿（不穿任何路径/mu线）
    \node[gray!55!black,anchor=west,font=\small] at (0.35,4.95)
        {波动带 $\mu\pm\dfrac{1.5\,\sigma}{\sqrt n}$};
    \draw[gray!55!black,->,thin] (2.45,4.78) -- (3.0,3.7);
\end{tikzpicture}
\caption{四条样本均值路径 $\{\bar X_n\}$（i.i.d.\ 抽样，$\mu=2$）。早期在宽处乱晃，随 $n$ 增大被一条按 $\propto 1/\sqrt n$ 收紧的波动带逼回水平线 $\mu$。强大数定律保证：除一个零概率的例外集外，\emph{每一条}这样的路径都收敛到 $\mu$。}
\label{fig:lln-2}
\end{figure}

\section{大数定律就是``一致性''}
\label{sec:lln-consistency}

把大数定律放进计量的语言，它的名字叫\textbf{一致性}（consistency）。一个估计量 $\hat\theta_n$ 称为 $\theta_0$ 的（弱）一致估计，如果 $\hat\theta_n\pto\theta_0$。对样本均值 $\bar X_n$ 估计 $\mu$ 这件事，``弱大数定律''与``$\bar X_n$ 是 $\mu$ 的一致估计''是同一句话的两种说法。

\state{无偏与一致是两回事，别混}{
\textbf{无偏}（unbiasedness）是\emph{固定样本量}下的性质：$\E{\hat\theta_n}=\theta_0$，平均而言瞄准靶心。\textbf{一致}（consistency）是\emph{样本量趋于无穷}下的性质：$\hat\theta_n\pto\theta_0$，估计量本身（而非其平均）越来越逼近真值。二者既不互推、也常常错位：样本均值\emph{既}无偏\emph{又}一致；样本方差除以 $n$（而非 $n-1$）的版本有偏，却仍一致；存在无偏却不一致的估计（如只用第一个观测 $\hat\theta_n=X_1$，永远无偏，但根本不收敛）。计量里真正不可或缺的是\emph{一致性}——它说``数据够多就对'',这恰是大数定律所给。
}

最常用、也最该记住的推广方式，是``连续函数不破坏一致性''。这一步把``样本均值''升级成``几乎任何由样本矩拼出来的估计量'',是从大数定律走向 GMM、MLE 一致性的关键转接口。

\thm{连续映射保持一致性}{
若 $\hat\theta_n\pto\theta_0$，且 $g$ 在 $\theta_0$ 处连续，则 $g(\hat\theta_n)\pto g(\theta_0)$。
}

这条结论的完整名字是\emph{连续映射定理}，将在后文``连续映射定理、Slutsky 与 Delta 法''一章里证明并推广（它对依分布收敛也成立）。它的威力在于：很多估计量本身不是某个简单平均，却是若干样本矩的\emph{连续函数}。比如样本方差
\[
    s_n^2=\frac1n\sum_{i=1}^n X_i^2-\bar X_n^2 ,
\]
它是两个样本矩 $\frac1n\sum X_i^2\pto\E{X^2}$ 与 $\bar X_n\pto\E{X}$ 的连续函数（$g(a,b)=a-b^2$）。于是由大数定律加连续映射，立刻有 $s_n^2\pto\E{X^2}-(\E{X})^2=\sigma^2$——样本方差是总体方差的一致估计。\textbf{大数定律给``一阶''的样本矩收敛，连续映射把它批发给一大类非线性估计量}。

\ex[样本相关系数的一致性]{
设 $(X_i,Y_i)$ i.i.d.，各阶矩有限。说明样本相关系数
\[
    \hat\rho_n=\frac{\frac1n\sum_i X_iY_i-\bar X_n\bar Y_n}{\sqrt{\pr{\frac1n\sum_i X_i^2-\bar X_n^2}\pr{\frac1n\sum_i Y_i^2-\bar Y_n^2}}}
\]
是总体相关系数 $\rho=\Corr(X,Y)$ 的一致估计。
}
\sol{
$\hat\rho_n$ 完全由五个样本矩搭成：$\frac1n\sum X_iY_i,\ \frac1n\sum X_i^2,\ \frac1n\sum Y_i^2,\ \bar X_n,\ \bar Y_n$。由大数定律，这五个样本矩分别依概率收敛到对应的总体矩
\[
    \E{XY},\quad \E{X^2},\quad \E{Y^2},\quad \E{X},\quad \E{Y} .
\]
而 $\hat\rho_n$ 是这五个量的一个函数 $g$，它在极限点处连续（只要分母里两个方差都为正，$g$ 在该点是初等运算与开方的复合，连续）。由连续映射定理，
\[
    \hat\rho_n\;\pto\;g\pr{\E{XY},\E{X^2},\E{Y^2},\E{X},\E{Y}}
    =\frac{\Cov{X}{Y}}{\sqrt{\Var{X}\Var{Y}}}=\rho .
\]
没有动用任何分布形式假设、也没碰收敛速度——一致性是大数定律加连续性``白送''的。
}

\section{它通向何处}
\label{sec:lln-applications}

大数定律不是概率论里一条孤立的极限定理，它是整张``用样本逼近总体''的网的中心结点。几处主要去处：

\begin{itemize}
    \item \textbf{矩估计与 GMM。} 矩估计直接拿样本矩当总体矩——其一致性\emph{就是}大数定律。GMM 把模型写成一组总体矩条件 $\E{g(\vc Z,\vc\theta_0)}=\vc 0$，估计量让对应的\emph{样本}矩 $\frac1n\sum_i g(\vc Z_i,\vc\theta)$ 尽量接近零；大数定律保证样本矩函数（逐点）收敛到总体矩函数，这是 GMM 一致性证明的第一块砖。
    \item \textbf{极大似然。} MLE 最大化样本平均对数似然 $\frac1n\sum_i\log f(X_i;\vc\theta)$。大数定律说它（逐点）收敛到其期望 $\E{\log f(X;\vc\theta)}$，而后者在真值处取最大（信息不等式）。``样本目标函数收敛到总体目标函数、各自的最大值点也随之靠拢''——这套论证（需配上一致收敛与可识别性）是 MLE 一致性的骨架，详见后文``极值估计的统一框架''一章。
    \item \textbf{蒙特卡洛积分。} 要算 $\E{h(X)}=\int h\,\d F$ 而积分写不出解析解时，抽 $X_1,\dots,X_n\iid F$、用 $\frac1n\sum_i h(X_i)$ 近似。这\emph{正是}大数定律：样本平均 $\pto$ 期望。整个随机模拟方法论（贝叶斯计算、结构估计里的模拟矩）都站在这一句上。
    \item \textbf{经验分布函数与 Glivenko--Cantelli。} 对每个固定的 $t$，经验分布 $\hat F_n(t)=\frac1n\sum_i\indicator_{\set{X_i\le t}}$ 是示性变量的样本均值，故 $\hat F_n(t)\pto F(t)$——又一次大数定律。Glivenko--Cantelli 定理把这一``逐点''收敛加强为\emph{一致}收敛 $\sup_t\abs{\hat F_n(t)-F(t)}\asto 0$，是``统计学基本定理'',为一切基于经验分布的方法（分位数估计、自助法）奠基。
    \item \textbf{时间序列的遍历定理。} 经济数据多是\emph{时间序列},观测前后相关，不满足 i.i.d.。此时把大数定律推广成\textbf{遍历定理}：对平稳遍历过程，时间平均 $\frac1n\sum_{t=1}^n X_t\asto\E{X_1}$ 仍成立。它让``用一条历史路径估计总体期望''在相依数据下重新合法，是宏观计量与金融计量的命脉。
\end{itemize}

\state{一句话收束}{
大数定律把``样本平均逼近总体期望''从信念变成定理；连续映射定理把这一条结论批发给几乎所有由样本矩搭成的估计量。\emph{一致性}——估计方法``数据够多就对''的最低承诺——其根全在这里。它的下一站是问``逼近得有多快、误差长什么样'',那就是中心极限定理的领地。
}
